\documentclass[10pt,a4paper]{amsart}
\usepackage[margin=19mm]{geometry}
\usepackage{amsmath,amssymb,mathtools}
\usepackage[scr=rsfs,cal=euler]{mathalfa}
\usepackage{xfrac}
\usepackage[unicode,hidelinks,bookmarks=false]{hyperref}
\newcommand{\ie}{i.e.\ }
\newcommand{\cf}{cf.\ }
\newcommand{\resp}{respectively\ }
\newcommand{\Subsection}{\subsection}
\providecommand{\nobreakdash}{\nobreak\hbox{-}}
\setlength{\emergencystretch}{2em}
\multlinegap0pt
\usepackage{amssymb}
\usepackage{tikz}
\usepackage{braket}
\usepackage{mathtools}
\usepackage{bm}
\usepackage{algorithm}\usepackage{algpseudocode}
\newtheorem{theorem}{Theorem}
\title{Finite-Precision Quantum Mechanics: Quantum Parcels, Information and Geometry}
\author{Abbas Edalat}
\date{Source: arXiv:2605.19706v5 (CC BY 4.0)}
\begin{document}
\maketitle

\noindent\emph{Source and licence.} Finite-Precision Quantum Mechanics: Quantum Parcels, Information and Geometry. Version arXiv:2605.19706v5 (CC BY 4.0), distributed by arXiv under the Creative Commons Attribution 4.0 International licence, http://creativecommons.org/licenses/by/4.0/, as recorded in the arXiv licence metadata for this exact version and shown on the abstract page https://arxiv.org/abs/2605.19706v5. Full licence notice: \texttt{licenses/arxiv-2605.19706-CC-BY-4.0.txt}.\par\smallskip

\noindent\textit{Editorial scope.} Counted unit: the article's theorem thm:volcontraction\_qubit (source0.tex lines 1270--1281). The article's complete original proof of it is delivered with it (source0.tex lines 2435--2443, one \texttt{proof} environment of 9 lines, which the article places away from the statement). No mathematics is written, repaired or re-derived: the statement and the proof are the article's own lines, in the article's own notation, and no retained line is substituted or re-flowed.  Retained uncounted as the source-local context the proof needs: lines 1293--1303.  Nothing was omitted from any retained span, so the package carries no omission record.  No retained line contains a citation, so the wrapper carries no bibliography and no bibliography entry.  No retained line contains a figure environment, an image-inclusion command or drawing code, so no image is delivered and the package declares no asset dependency.  Editorial formatting only, in addition: an AMS \texttt{amsart} preamble that restates the article's own declarations and macros used by the retained lines, namely the environments \texttt{theorem} on separate counters, together with the standard packages those lines need.  The retained environments print in this document as Theorem 1 in order of appearance, whereas the article prints the same items with the numbering of its own full document.  The complete native source of the article as distributed in the arXiv e-print for this exact version is bundled unchanged as \texttt{sources/200936/source0.tex}.
\begin{theorem}[Volume contraction, general finite-dimensional case]
\label{thm:volcontraction_nqubits}
Let \(P\subset\mathcal{D}(\mathcal{H})\) be a weak open subset of density matrices and assume
\[
\operatorname{Tr}(\rho\Pi_j)\ge c>0 \qquad\text{for every }\rho\in P.
\]
Then there exists \(\eta_0\in(0,1)\) such that for every \(\eta\in(\eta_0,1)\),
\[
\operatorname{Vol}_{\mathrm{HS}}(f_j^{(\eta)}(P))<\operatorname{Vol}_{\mathrm{HS}}(P).
\]
\end{theorem}

\begin{theorem}[Volume contraction for a qubit]
\label{thm:volcontraction_qubit}
Let \(P\subset\mathcal{D}(\mathbb{C}^2)\) be an open set of qubit states and assume that
\[
\operatorname{Tr}(\rho\Pi_j)\ge c>0 \qquad\text{for every }\rho\in P.
\]
Then there exists \(\eta_0\in(0,1)\) such that for every \(\eta\in(\eta_0,1)\),
\[
\operatorname{Vol}_{\mathrm{HS}}(f_j^{(\eta)}(P))<\operatorname{Vol}_{\mathrm{HS}}(P).
\]
If, moreover, \(c\ge\frac12\), then the inequality holds for all \(\eta\in(0,1)\).
\end{theorem}

\begin{proof}
Since $O_2\neq\varnothing$ is open, $\operatorname{Vol}_{\mathrm{HS}}(O_2)>0$. By construction $O_2\subseteq O_2'(\eta)$, so $\operatorname{Vol}_{\mathrm{HS}}(O_2'(\eta))\ge\operatorname{Vol}_{\mathrm{HS}}(O_2)>0$. By Theorem~\ref{thm:volcontraction_qubit} or Theorem~\ref{thm:volcontraction_nqubits}, $\operatorname{Vol}_{\mathrm{HS}}(O_1'(\eta))<\operatorname{Vol}_{\mathrm{HS}}(O_1)$. Therefore
\[
I_{\mathrm{HS}}(O_1'(\eta),O_2'(\eta))
=\frac{\operatorname{Vol}_{\mathrm{HS}}(O_2'(\eta))}{\operatorname{Vol}_{\mathrm{HS}}(O_1'(\eta))}
>\frac{\operatorname{Vol}_{\mathrm{HS}}(O_2)}{\operatorname{Vol}_{\mathrm{HS}}(O_1)}
=I_{\mathrm{HS}}(O_1,O_2).
\]
\end{proof}

\end{document}
